\section{Code and reproduction}
\label{app:code}

The code is in the repository \texttt{quantum-kicked-rotor} on GitHub. It uses Python 3 with NumPy, SciPy and
Matplotlib, and is organized as a small package plus one script per experiment:
\begin{itemize}
  \item \texttt{kicked\_rotor/classical.py}: the standard map and its inverse, the matched and torus
  ensembles, the Lyapunov exponent, and the exact classical two-kick result.
  \item \texttt{kicked\_rotor/quantum.py}: the momentum grid, the Floquet step and its inverse, and the exact
  quantum two-kick result.
  \item \texttt{kicked\_rotor/observables.py}: the norm, momentum distributions, $\avg{p^2}$, the edge
  probability, and the fits of the localization length.
  \item \texttt{kicked\_rotor/plotting.py}: shared figure style and saving.
\end{itemize}
Every function takes its parameters as arguments, with no hidden global state, so the package can be extended
for further numerical work.

\begin{table}[!htbp]
\centering
\small
\begin{tabular}{@{}llr@{}}
\toprule
Script & Figure & Run time \\
\midrule
\texttt{01\_classical\_phase\_portraits.py} & \cref{fig:phase_portraits} & 12 s \\
\texttt{02\_classical\_diffusion.py} & \cref{fig:diffusion} & 27 s \\
\texttt{03\_quantum\_vs\_classical.py} & \cref{fig:quantum_vs_classical} & 11 s \\
\texttt{04\_momentum\_distribution.py} & \cref{fig:distribution} and the animation & 50 s \\
\texttt{05\_resonance\_test.py} & \cref{fig:resonance} & 3 s \\
\texttt{06\_time\_reversal.py} & \cref{fig:time_reversal} & 6 s \\
\texttt{07\_bonus\_xi\_scaling.py} & \cref{fig:scaling} & 170 s \\
\bottomrule
\end{tabular}
\caption{The experiments in \texttt{experiments/}. Run times are from a Linux cloud machine; a recent laptop
is somewhat faster.}
\label{tab:experiments}
\end{table}

\paragraph{Running.}
On Windows, from the folder of a clone, with a virtual environment in \texttt{.venv}:
\begin{verbatim}
python -m venv .venv
.\.venv\Scripts\python -m pip install -r requirements.txt
.\.venv\Scripts\python experiments\03_quantum_vs_classical.py
.\.venv\Scripts\python -m pytest
\end{verbatim}
Each experiment prints a short summary and saves its figure to \texttt{figures/} as PNG and PDF. This report
is built from \texttt{report/} with \texttt{latexmk -pdf report.tex}.

\paragraph{Reproducibility.}
All random ensembles use fixed seeds, and rerunning any experiment on the same machine reproduces its figures
byte for byte. Across machines the quantum results agree to all printed digits, while the classical ones agree
only within their statistical errors, for the reason explained in \cref{sec:correspondence}. Experiment 07
sizes its runs from the five-term formula for $D$ rather than from the measured value, so that its quantum
results are machine independent too.
